% Folder 146, batch 01 (archivists' pages 1-20).
% Pass: Opus 5.5 (claude-opus-5-5), 2026-10-02 — first pass, unchecked against the pages by a human.
%
% Notation fixed for this batch: S_0\mathcal{T}_G (his "S_0 T_G", a script T
% with a lowered 0), \mathbb{U} for the circle group, P\mathbb{U}_{0,I} for the
% torsor of tangent directions, \mathfrak{S}_I for the symmetric group,
% \omega(I) for the set of the two orientations of I, \mathcal{E} for an
% extension group and \varepsilon for a torsor under its quotient.
\documentclass[11pt,a4paper]{article}
\input{../preamble/grothendieck}

\folder{146}
\batch{1}
\pages{1}{20}
\foldertitle{Groupoïdes de Teichmüller et les S0 Tg,!ν, S0 TG : notes manuscrites (s.d.).}
\dating{[à partir de 1978]}
\watermark{Édition de démonstration}

\begin{document}

\section{Calculs des $S_0\mathcal{T}_G$}
\note{titre souligné, de sa main, seul sur le feuillet 1, qui sert de couverture à la suite (feuillets 2 à 17)}

\page{2}
\subsection{I. Calcul de $S_0\mathcal{T}_G$, pour $G$ graphe MD de dim.\ modulaire $0$}

\marginal{\uncertain{Obtient}-on $S_0\mathcal{T}_G$ \ill{} $S_0\mathcal{T}_{G^b}$ par \ill{} \ill{} groupoïdes \ill{} \ill{} \ill{}}

$S_0\mathcal{T}_G$ est un $\mathbb{Z}^{\hat A}$-groupoïde, déduit du $\mathbb{Z}^{\vec{\hat A}}$-groupoïde $S_0\mathcal{T}_{G^b}$ par le changement de groupes \struck{\ill{}} interne
\[
  \mathbb{Z}^{\vec{\hat A}} \longrightarrow \mathbb{Z}^{\hat A}
\]
\uncertain{associé} à : $\vec{\hat A} \longrightarrow \hat A$,
i.e.
\[
  (1) \qquad S_0\mathcal{T}_G = S_0\mathcal{T}_{G^b} \wedge^{\mathbb{Z}^{\vec{\hat A}}} \mathbb{Z}^{\hat A}
\]
et
\[
  (2) \qquad S_0\mathcal{T}_{G^b} = \prod_{s\in S} S_0\mathcal{T}_{G^b_s} = \prod_{s\in S} S_0\mathcal{T}_{0,\vec{\hat A}_s} ,
\]
ce qui nous ramène à : \uncertain{décrire} $S_0\mathcal{T}_G$, quand $G$ est un \struck{\ill{}}graphe ayant un seul sommet, i.e.\ $G$ est un $G_{0,I}$, $I$ étant un ens.\ à trois éléments, \struck{de sorte} --- i.e.\ on est ramené à décrire \struck{\ill{}} $S_0\mathcal{T}_{0,I}$.

Considérons la \struck{courbe} \add{droite} projective $I$-marquée sur $\mathbb{C}$, $\mathbb{P}^1_I$. Soit $\omega(I)$ l'une des deux orientations de $I$. Si $\alpha \in \omega(I)$, on a des iso.\ canoniques (par les espaces tangents) \add{de $\mathbb{U}$-torseurs} \ill{} $T_i$ aux $i \in I \subset \mathbb{P}^1_I$
\[
  T_i \xleftrightarrow[\ \sim\ ]{\ \varphi_{\alpha,i}\ } \mathbb{U}
  \qquad \mathbb{U} = \{ z \in \mathbb{C} \mid |z| = 1 \}
\]
i.e.\ un isom.\ de $\mathbb{U}^I$-torseurs
\[
  \mathbb{U}^I \xrightarrow[\ \sim\ ]{\ \varphi\ } P\mathbb{U}_G = P\mathbb{U}_{0,I} = \prod_{i\in I} T_i
\]
i.e.\ un point \struck{noté} $\bar\alpha = \varphi(1) \in P\mathbb{U}_{0,I}$, d'où une \struck{\ill{}}\uncertain{injection} canonique
\[
  (3) \qquad \omega(I) \hookrightarrow P\mathbb{U}_{0,I}
\]
compatible avec l'homomorphisme sur les groupes d'opérateurs

\note{diagramme numéroté (4) dans la marge}

\begin{tikzcd}
\{\pm 1\} \arrow[rr] \arrow[dr, "\mathrm{diag}"'] & & \mathbb{U}^I \\
& \{\pm 1\}^I \arrow[ur, "\text{déduit de l'inclusion } \{\pm1\} \hookrightarrow \mathbb{U}"'] &
\end{tikzcd}

\page{3}
Bien entendu, (3) et (4) sont compatibles avec l'action de $\mathfrak{S}_I$ (qui opère sur $\omega(I)$ par la signature).
Je définis $S_0\mathcal{T}_{0,I}$ comme le \uncertain{sous}-groupoïde du groupoïde fondamental de $P\mathbb{U}_{0,I}$ induit par son ens.\ d'objets $\omega(I)$. Je voudrais en donner une expression purement algébrique. Pour ceci, soit $\tilde{\mathbb{U}} \simeq \mathbb{R}$ le \uncertain{revêtement} universel standard de $\mathbb{U}$, et considérons la suite exacte
\[
  (5) \qquad 0 \to \mathbb{Z} \to \mathbb{R} \to \mathbb{U} \to 0 ,
  \qquad t \mapsto \exp 2i\pi t
\]
d'où
\[
  (5_I) \qquad 0 \to \mathbb{Z}^I \to \mathbb{R}^I \to \mathbb{U}^I \to 0 ,
\]
Notons que $P\mathbb{U}_{0,I}$ est un torseur sous le groupe topologique $\mathbb{U}^I$, \struck{\ill{}} \add{structure algébrique} qui détermine sa topologie. Je vais \ill{} une construction purement algébrique \uncertain{du} groupoïde fondamental \ill{} \ill{}, en utilisant seulement $(5_I)$ et sa structure de torseur. Plus généralement, considérons une structure d'extension de groupes topologiques
\[
  (6) \qquad 0 \to \pi \to \tilde T \to T \to 0
\]
$\pi$ discret, avec $\pi \subset \tilde T^0$ et $\tilde T^0$ simplement connexe

\page{4}
(\struck{contractile}\add{simplement connexe}), $\tilde T$ est un espace contractile tel que \ill{} de rang fini, $\pi$ un $\mathbb{Z}$-module libre de t.f.).
\struck{On sait, pour $x, y \in$} \add{Soit $P$ un $\mathbb{U}$-torseur,} qui hérite donc d'une topologie, d'où \ill{} un groupoïde fondamental ? \uncertain{Peut}-on décrire purement algébriquement le groupoïde \ill{} ? \ill{} on \ill{}, pour $x, y \in P$, on a \ill{}
\[
  (7) \qquad \Pi(x,y) = \tilde{\mathbb{U}}_{y-x}
  \qquad \text{(fibre de $\tilde{\mathbb{U}}$ en $y - x \in \mathbb{U}$)}
\]
La composition des flèches dans $\Pi$ \ill{} étant donnée par l'addition dans $\tilde{\mathbb{U}}$.

\struck{\ill{}} Cette formule décrit \struck{un groupoïde} en fait, pour toute extension de groupes abéliens
\[
  (8) \qquad 0 \to \pi \to \mathcal{E} \to \Gamma \to 0
\]
et un torseur $\varepsilon$ sous $\Gamma$ --- on \add{connaît} un $\pi$-groupoïde $\mathcal{E}(\varepsilon)$, ayant $\varepsilon$ comme ens.\ d'objets, et ses flèches données par
\[
  (9) \qquad \mathrm{Hom}(x,y) = \mathcal{E}_{y-x}
\]
la composition des flèches étant donnée par l'addition dans $\mathcal{E}$. On peut aussi décrire \uncertain{ce} groupoïde par l'ens.\ des flèches
\[
  (10) \qquad \mathrm{Fl}\,\mathcal{E}(\varepsilon) = \mathcal{E} \times \varepsilon
\]
\note{un symbole biffé, illisible, devant $\mathcal{E} \times \varepsilon$}

\marginal{\ill{} \ill{} un $\pi$-groupoïde \ill{} \ill{} connexe \ill{} \ill{} \ill{} $\mathcal{E}(\varepsilon)$ \ill{} \ill{} (9) \ill{}}

\page{5}
Avec l'application source donnée par $\mathrm{pr}_2$, l'application but par la loi d'opération de $\mathcal{E}$ sur $\varepsilon$ (via $\Gamma$ et $\mathcal{E} \to \Gamma$), la loi de composition étant donnée par
\[
  (u, vx) \circ (v, x) = (uv, x)
\]
Considérons un groupe $\Gamma'$, un \uncertain{$\Gamma'$}-torseur $\varepsilon'$, et un homomorphisme de torseurs
\[
  \varepsilon' \longrightarrow \varepsilon
\]
associé à un homomorphisme de groupes
\[
  \Gamma' \longrightarrow \Gamma
\]
On veut décrire le $\pi$-groupoïde induit sur $\varepsilon'$ comme ens.\ d'objets par le $\pi$-groupoïde $\mathcal{E}(\varepsilon)$, \ill{} : l'application \ill{}. On trouve qu'il est canon.\ isomorphe au $\pi$-groupoïde $\mathcal{E}'(\varepsilon')$, où $\mathcal{E}'$
\[
  (11) \qquad 0 \to \pi \to \mathcal{E}' \to \Gamma' \to 0
\]
est l'extension, image inverse de (8) par $\Gamma' \to \Gamma$.

Remarque : la description de $S_0\mathcal{T}_{0,I}$, en termes de l'inclusion (3), se traduit :
\[
  (12) \qquad S_0\mathcal{T}_{0,I} \simeq \mathcal{E}_I(\omega(I))
\]
\add{comme} $\mathbb{Z}^I$-groupoïde,
où $\mathcal{E}_I$ est l'extension
\[
  (13) \qquad 0 \to \mathbb{Z}^I \to \mathcal{E}_I \to \{\pm 1\} \to 0
\]

\page{6}
image inverse de l'extension $(5_I)$, par l'homomorphisme (4) de $\underbrace{\{\pm1\}}_{\mu_2}$ dans $\mathbb{U}^I$.

Je vais donner une description plus simple de cette extension, en \struck{\ill{}} \add{factorisant} \struck{\ill{}} \ill{} de (4)
\[
  \mu_2 \xrightarrow{\ \mathrm{diag}\ } \mu_2^I \longrightarrow \mathbb{U}^I .
\]
\uncertain{D'abord} \ill{} \ill{} : déterminons l'ext.\ \ill{} de $\mu_2$ par $\mathbb{Z}$ déduite de l'ext.\ can.\ (5) de $\mathbb{U}$ par $\mathbb{Z}$, par le \uncertain{plongement} $\mu_2 \subset \mathbb{U}$ --- extension qui est canoniquement isomorphe à
\[
  (14) \qquad 0 \to \mathbb{Z} \xrightarrow{\ 2\ } \mathbb{Z} \to \mathbb{Z}/2\mathbb{Z} \to 0 ,
  \qquad \mathbb{Z}/2\mathbb{Z} \simeq \mu_2 .
\]
\uncertain{De là} on prend le produit cartésien $I$-ième \ill{}
\[
  0 \to \mathbb{Z}^I \xrightarrow{\ 2\ } \mathbb{Z}^I \to (\mathbb{Z}/2\mathbb{Z})^I \to 0 ,
  \qquad (\mathbb{Z}/2\mathbb{Z})^I \simeq \mu_2^I ,
\]
qui donne l'image inverse de $(5_I)$ par $\mu_2^I \to \mathbb{U}^I$, puis on prend l'ext.\ induite (de $\mu_2$ par $\mathbb{Z}^I$) déduite de l'homom.\ diagonal $\mu_2 \to \mu_2^I$. On peut considérer l'homom.\ d'extensions

\begin{tikzcd}
0 \arrow[r] & \mathbb{Z} \arrow[r, "2"] \arrow[d, "\mathrm{diag}"'] & \mathbb{Z} \arrow[r] \arrow[d, "\mathrm{diag}"] & \mathbb{Z}/2\mathbb{Z} \arrow[r] \arrow[d, "\mathrm{diag}"] & 0 \\
0 \arrow[r] & \mathbb{Z}^I \arrow[r, "2"] & \mathbb{Z}^I \arrow[r] & (\mathbb{Z}/2\mathbb{Z})^I \arrow[r] & 0
\end{tikzcd}

\page{7}
qui donne la description suivante de (13) comme extension déduite de l'extension évidente de $\mathbb{Z}/2\mathbb{Z}$ par $\mathbb{Z}$, grâce à l'hom.\ diag.\ $\mathbb{Z} \to \mathbb{Z}^I$ :

\note{diagramme numéroté (15)}

\begin{tikzcd}
0 \arrow[r] & \mathbb{Z} \arrow[r, "2"] \arrow[d] & \mathbb{Z} \arrow[r] \arrow[d] & \mathbb{Z}/2\mathbb{Z} \arrow[r] \arrow[d, "\wr"] & 0 \\
0 \arrow[r] & \mathbb{Z}^I \arrow[r] & \mathcal{E}_I \arrow[r] & \mu_2 \arrow[r] & 0
\end{tikzcd}

Donc pour un $\mu_2$-torseur $\varepsilon$ i.e.\ un ens.\ à deux éléments $\varepsilon$, on a pour $x, y \in \varepsilon$ un $\mathbb{Z}^I$-torseur
\[
  (16) \qquad \underbrace{\mathrm{Hom}}_{\mathcal{E}_I(\varepsilon)}(x,y)
  = T_{x-y} \wedge^{\mathbb{Z}} \mathbb{Z}^I \simeq
  \begin{cases}
    \mathbb{Z}^I & \text{si } x = y \\
    \mathrm{Imp}^I & \text{si } x \neq y
  \end{cases}
\]
\marginal{où $\mathrm{Imp}$ = l'ens.\ des entiers impairs.}

où \struck{$T_{x-y}$} $T_{x-y}$ désigne le $\mathbb{Z}$-torseur \add{des} image inverse de $x - y \in \mathbb{Z}/2\mathbb{Z}$ dans $\mathbb{Z}$ (i.e.\ formé des entiers pairs si $x = y$, des entiers impairs si $x \neq y$), et où l'extension du groupe d'opérateurs de torseurs est \uncertain{celle} de : l'hom.\ diagonal $\mathbb{Z} \to \mathbb{Z}^I$.

Ceci s'applique en particulier au cas de $\mathcal{E}_I(\omega(I))$ i.e.\ où $\varepsilon = \omega(I)$, et donne la description algébrique complète de $S_0\mathcal{T}_{0,I}$, en termes de $I \in \mathrm{Ob}\,\mathrm{Ens}$.

Cette description est compatible avec les opérations de $\mathfrak{S}_I$. On trouve donc une

\page{8}
description des groupes fondamentaux \uncertain{variants} en un point $x \in \omega(I)$ :
\[
  (17) \qquad \pi_1(S_0\mathcal{T}_{0,I}, \mathfrak{S}_I; x)
  = \bigl\{ \alpha, l \bigm| \alpha \in \mathfrak{S}_I,\ l : x \leftarrow \alpha x \text{ dans } S_0\mathcal{T}_{0,I} \bigr\}
\]
\[
  = \bigl\{ (\alpha, (n_i)_{i\in I}) \in \mathfrak{S}_I \times \mathbb{Z}^I \bigm|
  \text{les $n_i$ ont la parité définie par $\mathrm{sg}(\alpha)$} \bigr\}
\]
\note{au-dessus de (17), un premier essai biffé : \struck{$S_0\mathcal{T}_{0,I}$}}
avec \uncertain{composition}
\[
  (18) \qquad (l, \alpha)(l', \alpha') = (\underbrace{l \circ \alpha(l')}_{l + \alpha(l')}, \alpha\alpha')
\]
\note{dans (18), un mot biffé, illisible, devant $l \circ \alpha(l')$}
\marginal{De façon générale, si $x, y \in \varepsilon$ ($= \omega(I)$), les hom.\ dans le $\Pi_1$ \uncertain{variant} de $x$ dans $y$ sont donnés par les couples $(l, \alpha)$ avec $l : \alpha x \to y$, et \ill{} \ill{} \ill{} par (18)}

Comme $S_0\mathcal{T}_{0,I}$ est déduit d'un $\mathbb{Z}$-groupoïde (sur $\omega(I)$ comme ens.\ d'objets) par changement des groupes d'opérateurs \ill{} $\mathbb{Z} \to \mathbb{Z}^I$, l'\uncertain{extension} \ill{} de $\mathfrak{S}_I$ par $\mathbb{Z}^I$ peut être considérée comme \ill{} déduite d'une extension de $\mathfrak{S}_I$ par $\mathbb{Z}$ (savoir le groupe fond.\ variant \ill{} précédent) \ill{} \ill{} \ill{} $\mathbb{Z}$-groupoïdes \ill{} (à opérateurs $\mathfrak{S}_I$).
Sa description explicite est \ill{} \ill{} \struck{\ill{}} \ill{} (17), (18), \ill{} que les $l$ sont des éléments de $\mathbb{Z}$ \ill{} lieu de $\mathbb{Z}^I$ --- et l'op.\ de $\mathfrak{S}_I$ sur $\mathbb{Z}$ est triviale. Donc on voit que l'ext.\ \ill{} \ill{} \ill{} que l'extension, \uncertain{déduite} \ill{} de l'extension (14) de $\mu_2$ par $\mathbb{Z}$, \ill{} \ill{} l'homom.\ signature $\mathfrak{S}_I \to \mu_2$. En résumé :

\textit{Proposition.} --- Soit $x \in \omega(I)$. Alors $\pi_1(S_0\mathcal{T}_{0,I}, \mathfrak{S}_I; x)$ est can.\ isom.\ à \uncertain{une} l'extension de $\mathfrak{S}_I$ par $\mathbb{Z}^I$

\page{9}
obtenue à partir de l'extension (14) de $\mathbb{Z}/2\mathbb{Z} \simeq \mu_2$ par $\mathbb{Z}$
\[
  0 \to \mathbb{Z} \xrightarrow{\ 2\ } \mathbb{Z} \to \underset{\mu_2}{\mathbb{Z}/2\mathbb{Z}} \to 0 ,
\]
en prenant l'image inverse de celle-ci par l'homom.\ signature
\[
  \mathfrak{S}_I \xrightarrow{\ \mathrm{sgn}\ } \mu_2 \simeq \mathbb{Z}/2\mathbb{Z} ,
\]
et en prenant la $\mathbb{Z}^I$-extension déduite de celle-ci par l'homom.\ \add{diag.} $\mathbb{Z} \to \mathbb{Z}^I$.

Cette extension n'est pas scindable --- ou l'extension induite sur un sous-groupe \add{$\Gamma_k$} d'ordre $2$, défini par une transposition $\sigma_k = (i,j)$, ne l'est pas --- ou en \uncertain{augmentant} la projection $\mathbb{Z}^I \to \mathbb{Z}$ correspondant aux \ill{} indices \ill{} \ill{} par l'\uncertain{involution}) --- on trouve une extension de $\Gamma_k$ par $\mathbb{Z}$, qui n'est autre (via l'identification $\Gamma_k \simeq \mathbb{Z}/2\mathbb{Z}$) que l'extension (14) de $\mathbb{Z}/2\mathbb{Z}$ par $\mathbb{Z}$, qui n'est pas scindable évidemment. Par contre la restriction de l'extension à $\mathfrak{S}_I^+$, qui est le \uncertain{stabilisateur} de $x$ dans $\mathfrak{S}_I$, est évidemment canoniquement scindée --- ce qui est manifeste aussi sur la formule explicite (17), (18) --- le scindage est
\[
  (19) \qquad \alpha \longmapsto (0, \alpha) \qquad \alpha \in \mathfrak{S}_I^+ .
\]

\page{10}
Il est \uncertain{un peu étrange} que l'expression obtenue pour $\pi_1(S_0\mathcal{T}_{0,I}, \mathfrak{S}_I; x)$ ne dépende pas du choix de $x \in \omega(I)$ --- c'est un groupe $G(I)$, extension de $\mathfrak{S}_I$ par $\mathbb{Z}^I$, qui se construit canoniquement, sans référence à un $x \in \omega(I)$. Donc $\mathfrak{S}_I$ opère \ill{} de façon compatible avec son \ill{} (\ill{}), bien qu'il n'y ait pas de scindage par quoi décrire cette opération. Bien entendu, la restriction de cette opération à $\mathfrak{S}_I^+$ n'est autre que celle déduite du scindage canonique au-dessus de $\mathfrak{S}_I^+$.

Regardons maintenant le cas d'un graphe MD\struck{G} $G$ de dim.\ modulaire $0$, de sorte que $S_0\mathcal{T}_G$ se déduit du produit
\[
  S_0\mathcal{T}_{G^b} = \prod_{s\in S} S_0\mathcal{T}_{G^b_s} = \prod_{s\in S} S_0\mathcal{T}_{0,\vec{\hat A}_s}
\]
par l'homom.\ des groupes \struck{\ill{}} d'op.\ internes
\[
  (20) \qquad \mathbb{Z}^{\vec{\hat A}} \longrightarrow \mathbb{Z}^{\hat A} .
\]
On peut d'autre part décrire le produit $S_0\mathcal{T}_{G^b}$ en termes des torseurs
\[
  \varepsilon(G) = \varepsilon(G^b) = \prod_{s\in S} \underbrace{\omega(\vec{\hat A}_s)}_{\varepsilon(G^b_s)}
\]
\note{la ligne se poursuit au feuillet suivant ; en marge, devant elle, un mot abrégé, « \uncertain{et} »}

\page{11}
\struck{i.e.\ l'ext.\ \ill{}} sous le groupe
\[
  (22) \qquad \gamma(G) = \gamma(G^b) = \{\pm 1\}^S
\]
et de l'extension
\[
  (23) \qquad 0 \to \mathbb{Z}^S \xrightarrow{\ 2\ } \mathbb{Z}^S \to \underset{\gamma(G)}{(\mathbb{Z}/2\mathbb{Z})^S} \to 0
\]
(puiss.\ cartésienne, d'exposant $S$, de l'ext.\ (14)) en considérant l'homomorphisme
\[
  \mathbb{Z}^S \longrightarrow \mathbb{Z}^{\vec{\hat A}}
\]
déduit de l'application-source
\[
  \vec{\hat A} \longrightarrow S
\]
d'où une extension
\[
  (24) \qquad 0 \to \mathbb{Z}^{\vec{\hat A}} \longrightarrow \mathcal{E}(G^b) \longrightarrow \gamma(G) \to 0
\]
(ne dépendant que de $G^b$). Ceci posé, on a
\[
  (25) \qquad S_0\mathcal{T}_{G^b} = \mathcal{E}(G^b)(\underbrace{\varepsilon(G^b)}_{\varepsilon(G)})
\]
Pour passer à $S_0\mathcal{T}_G$, on considère l'extension
\[
  (26) \qquad 0 \to \mathbb{Z}^{\hat A} \longrightarrow \mathcal{E}(G) \longrightarrow \gamma(G) \to 0
\]
de $\gamma(G) = \gamma(G^b)$ par $\mathbb{Z}^{\hat A}$, déduite de $\mathcal{E}(G^b)$ (24) par l'hom.\ (20). On trouve
\note{les numéros d'équations (22) à (26) de ce feuillet sont lus avec peine, plusieurs chiffres étant récrits ; la suite (26) est reprise au feuillet suivant}

\page{12}
l'expression intrinsèque
\[
  (27) \qquad S_0\mathcal{T}_G \simeq \mathcal{E}(G)(\varepsilon(G))
\]
i.e.
\[
  (28) \qquad
  \begin{cases}
    \mathrm{Ob}\, S_0\mathcal{T}_G = \varepsilon(G) \\
    \mathrm{Hom}_{S_0\mathcal{T}_G}(x,y) = \underbrace{\mathcal{E}(G)_{y-x}}_{\text{fibre de $\mathcal{E}(G)$ au-dessus de $y - x \in \gamma(G)$}}
    \quad \text{pour } x, y \in \varepsilon(G)
  \end{cases}
\]
la composition des homomorphismes étant donnée par la \struck{\ill{}} \add{loi d'addition} dans $\mathcal{E}(G)$.

\marginal{On pourrait d'ailleurs prendre $\mathcal{E}(G)$ \uncertain{et} $\varepsilon(G)$ \ill{} (27) \ill{} (23), \ill{} $\mathbb{Z}^{\hat A} \leftarrow \mathbb{Z}^{\vec{\hat A}} \leftarrow \mathbb{Z}^S$}

On peut prendre (28) comme \textit{définition} de $S_0\mathcal{T}_G$, pour $G$ de dim.\ modulaire nulle.

Considérons p.ex.\ le \struck{\ill{}} \add{MD-graphe} du type

\drawing{un arbre à deux sommets reliés par l'arête $a$, chaque sommet portant en outre deux demi-arêtes libres, terminées par des flèches}

correspondant : à la donnée d'un ens.\ \struck{\ill{}} \add{$I$} de cardinal $4$, et d'une partition de type $(2,2)$ de $I$, i.e.\ à la donnée d'une \textit{courbe} combinatoire \uncertain{à} \ill{} $I$ \ill{} l'une des \uncertain{sommets}, et la partition \ill{} \ill{} \uncertain{celle} des \ill{} \ill{}). On \uncertain{obtient}

\page{13}
$G_Q$ le graphe associé à un carré combinatoire $Q$,

\drawing{un rectangle, les quatre sommets marqués d'un point}

Soit $\mathrm{diag}\,Q$ l'ens.\ des deux diagonales de $Q$, \struck{$A(Q)$} $A(Q)$ l'ens.\ des quatre sommets, s'identifiant à
\[
  A(Q) = \prod_{\Delta \in \mathrm{diag}\,Q} \Delta
\]
\struck{ens.\ sous-jacent} \add{torseur sous} $\{\pm1\}^{\mathrm{diag}\,Q}$, \struck{qui}
\struck{\ill{}} \ill{} : $\{\pm1\} \subset \{\pm1\}^{\mathrm{diag}\,Q}$ \add{y opère} par passage du sommet opposé, \struck{\ill{}} \uncertain{à} \uncertain{codiagonales}. \uncertain{Donc}
\[
  \mathrm{codiag}(Q) \simeq \bigwedge_{\Delta \in \mathrm{diag}\,Q} \Delta .
\]
On a ici, en termes du graphe $G_Q$
\[
  \begin{aligned}
    S &\simeq \mathrm{diag}\,Q && \text{soit $\Delta_s$ la diagonale associée à $s \in S$} \\
    I = A^{\mathrm{lib}} &= A(Q) && \text{ens.\ des sommets de $Q$} \\
    \vec{\hat A} &\simeq A \amalg \tilde S && \text{en associant à $s \in S$ l'arête orientée $\tilde s$ d'origine $s$}
  \end{aligned}
\]
\note{dans la deuxième ligne, un symbole biffé devant $A(Q)$, et « sommets » écrit au-dessus d'un mot biffé}
\[
  \vec{\hat A}_s = \Delta_s \amalg \{\tilde s\}
  \qquad \text{donc} \qquad
  \varepsilon_s = \omega(\vec{\hat A}_s) \simeq \Delta_s
\]
\[
  \varepsilon(G) \simeq \prod_{s\in S} \Delta_s \simeq A(Q)
  \qquad \text{en tant que torseurs sous}
\]
\[
  \gamma_G = \prod_{s\in S} \{\pm1\} = \{\pm1\}^S .
\]
On trouve $S_0\mathcal{T}_G$ à l'aide de l'extension $\mathcal{E}_{G^b}$
\[
  0 \longrightarrow \mathbb{Z}^S \xrightarrow{\ 2\ } \mathbb{Z}^S \longrightarrow \overbrace{\{\pm1\}^S}^{\gamma(G)} \longrightarrow 0
\]
et du \add{$\gamma(G)$-torseur} $\varepsilon(G)$, en prenant $\mathcal{E}_{G^b}(\varepsilon(G))$, \add{sur $\gamma(G)$,} qui donne un $\mathbb{Z}^S$-\struck{groupoïde}, et en \uncertain{faisant} l'extension des groupes internes $\mathbb{Z}^S \to \mathbb{Z}^{\vec{\hat A}}$, puis

\page{14}
l'homomorphisme

\begin{tikzcd}[column sep=small]
\mathbb{Z}^S \arrow[rr] \arrow[dr, "\text{déduit par contravariance de } \vec{\hat A} \xrightarrow{\mathrm{or.}} S"'] & & \mathbb{Z}^{\hat A} \\
& \mathbb{Z}^{\vec{\hat A}} \arrow[ur, "\text{déduit par cov.\ de } \vec{\hat A} \to \hat A \simeq \vec{\hat A}/\sigma"'] &
\end{tikzcd}

\note{à droite du diagramme, une ligne biffée qui se termine par $\vec{\hat A} \to S$, $\vec a \mapsto \mathrm{or}(\vec a)$ : \struck{\ill{}}}

On a : $\mathbb{Z}^{\hat A} \simeq \mathbb{Z}^I \times \mathbb{Z}$, et $\mathbb{Z}^S \to \mathbb{Z}^I \times \mathbb{Z}$ se décompose en
\[
  \mathbb{Z}^S \to \mathbb{Z}^I \quad \text{associé par contravariance à } I \xrightarrow[\deg 2]{} S
\]
et en
\[
  \mathbb{Z}^S \xrightarrow[\text{homm.\ somme}]{} \mathbb{Z} .
\]

\page{15}
\subsection{II. Détermination de $S_0\mathcal{T}_{0,I}$, pour $\mathrm{card}\, I = 4$}

Soit
\[
  (1) \qquad
  \begin{cases}
    J = \text{ens.\ des partitions de type $(2,2)$ de $I$} \\
    \Sigma = \mathbb{P}^1_J \quad \text{la \struck{sphère} \add{dr.\ proj.} $J$-marquée associée} \\
    \omega = \omega(J) \quad \text{l'ens.\ des orientations de $J$.}
  \end{cases}
\]
\marginal{N.B.\ $\Sigma = \Sigma_I$ \struck{\ill{}} \uncertain{opère} sur $\mathfrak{S}_I$ par l'intermédiaire de $\mathfrak{S}_I \to \mathfrak{S}_J$, \uncertain{de noyau} $V(I) \simeq$ \ill{} (le groupe de Klein) \ill{} \ill{} \ill{} \ill{} \ill{} $\mathbb{F}_2$, \uncertain{déterminé} par $J$ : $V(I)^* \simeq J$}

On peut plonger $\omega$ dans $\Sigma$ par un isomorphisme
\[
  (2) \qquad \omega \simeq \Sigma^{\mathfrak{S}_J^+} = \Sigma^{\mathfrak{S}_I^+}
\]
en associant : $x \in \omega$ le point fixe sous $\mathfrak{S}_J^+$ qui est dans l'hémisphère qui correspond : l'orientation de l'équateur $\Sigma_{\mathbb{R}}$ correspondante : l'ordre circulaire donné sur $J$. Pour \struck{\ill{}} $\alpha \in J$, je désigne par \struck{$\sigma_\alpha$ la transposition} $\sigma_\alpha \in \mathfrak{S}_J$ la transposition dans $J$ qui fixe $\alpha$, \struck{$\tilde\alpha$}\add{$\tilde\alpha$} le pt fixe \struck{\ill{}} $\neq \alpha$ de cette transposition.

\drawing{une sphère : pôle $P$ en haut et pôle opposé en bas ; l'équateur porte les points $\beta$, $\bar\alpha$ et un troisième point au bord droit ; un parallèle tracé en pointillé porte $\tilde\alpha$, $\alpha$, $\tilde\beta$}
\marginal{la diagonale \ill{} \ill{} est \ill{} \ill{} \ill{} \ill{} $\bar\Sigma$ par $P$ et \ill{}}

On pose
\[
  (3) \qquad \Sigma^* = \Sigma_I^* = \Sigma \setminus J
\]
\note{deux symboles biffés et illisibles, l'un après $\Sigma^*$, l'autre en exposant de $\Sigma_I$, au-dessus duquel on lit un $\alpha$}
On a des iso.\ canoniques (sur $\mathbb{C}$)
\[
  (4) \qquad
  \begin{cases}
    M_{0,I} \simeq \Sigma_I^* \\
    \hat M_{0,I} \simeq \Sigma_I
  \end{cases}
\]

\page{16}
On prend
\[
  \Sigma_J = \mathbb{P}(V(J)) \simeq \check{\mathbb{P}}(V(J))
\]
où $V(J)$ défini par une \uncertain{suite exacte}
\[
  0 \to V(J) \longrightarrow \mathcal{O}^J \xrightarrow[\text{somme}]{} \mathcal{O} \to 0
\]
La section $\alpha^*$ de $\Sigma_J$ est définie par la \uncertain{suite}

\struck{la section \ill{} $\Sigma$ \ill{} $\beta$ de $V($}
\[
  V(J) \xrightarrow{\ \mathrm{pr}_\alpha\ } \mathcal{O}
\]
\note{sous la ligne biffée, également biffé : « pt de $\{\alpha\}$ » ; à droite de la flèche, deux expressions biffées $\mathcal{O}.(\beta - \gamma)$, \ill{}, où $\beta$, $\gamma$ sont les \ill{}}
qui est donnée \struck{par} \add{loc.} par \uncertain{les} \uncertain{points} \ill{} $J$ \ill{} de $\alpha$. Quand : $\alpha'$, il est donné par la section
\[
  -2\alpha + \beta + \gamma = -2\alpha + \sum_{\beta \in J \setminus \{\alpha\}} \beta
\]
soit $\Delta_{\alpha'}$. On a
\[
  0 \to \Delta_{\alpha'} \longrightarrow V(J) \longrightarrow \mathcal{O}_{\alpha'}(1) \to 0
\]
\uncertain{d'où} \add{et}
\[
  \Delta_{\alpha'} \xrightarrow[\substack{\mathrm{pr}_\beta \\ (\beta \neq \alpha)}]{\ \sim\ } \mathcal{O} .
\]
isom.\ indépendant du choix de $\beta$.
\uncertain{On} trouve donc, prenant les déterminants
\[
  \det V(J) \simeq \mathcal{O}_{\alpha'}(1) \simeq \mathcal{O}(1) \otimes_{\mathbb{Z}} \mathbb{Z}(\omega)
\]
\note{le second isomorphisme est écrit verticalement sous le premier ; $\mathbb{Z}(\omega)$ y est glosé : « $\mathbb{Z}$ tordu par $\omega = \omega(J)$ »}
On a un isom.
\[
  \hat T_i \simeq \mathcal{O}(1)^{\otimes} \otimes_{\mathbb{Z}} L(\omega) ,
\]
\note{deux symboles biffés, illisibles, devant $\hat T_i$ et devant $\mathcal{O}(1)$ ; l'exposant de $\mathcal{O}(1)$ est incertain}
où $L(\omega)$ est l'alternative
\[
  \begin{cases}
    L(\omega) = \mathbb{Z} \\
    L(\omega) = \mathbb{Z}(\omega)
  \end{cases}
\]
\note{en regard de la première ligne : « (\uncertain{non} \ill{} \ill{}, $i$) » ; de la seconde : « (id.) où $\omega = \omega(J)$ »}
\marginal{\uncertain{Signification} de $\hat T_{0,I}$ dans $\hat M_{0,I}$ : le \ill{} $i \in I$ \ill{} de la section \ill{}}

Pour déterminer quelle est l'alternative correcte

\page{17}
Je prends l'image inverse par la section \struck{\ill{}} $\alpha'$ ($\alpha \in J$) dans $M_{0,I}$, on aura donc
\[
  \underset{\substack{\text{fibré tangent} \\ \text{en long de la} \\ \text{section $i$ de la} \\ \text{droite proj.\ relative} \\ \text{$X$ $I$-marquée}}}{T_i}
  \simeq \underbrace{\underbrace{\mathcal{O}_{\alpha'}(1)}_{\mathcal{O} \otimes_{\mathbb{Z}} \mathbb{Z}(\omega)} \otimes_{\mathbb{Z}} L(\omega)}_{\mathcal{O} \otimes_{\mathbb{Z}} (\mathbb{Z}(\omega) \otimes_{\mathbb{Z}} L(\omega))}
\]
\note{devant $\mathcal{O}_{\alpha'}(1)$, un symbole biffé}
Mais avec $\alpha$, pour une structure carrée $Q = Q_\alpha$ sur $I$, et un sommet marqué $i$, d'où une diagonale marquée

\drawing{un carré, le sommet en bas à gauche marqué $i$}

\[
  \omega(Q) \simeq \omega
\]
\note{après $\omega$, un symbole biffé et noirci}
et on trouve
\[
  T_i \simeq \mathcal{O} \otimes_{\mathbb{Z}} \mathbb{Z}(\omega)
\]
et comparant, on trouve $L(\omega) = \mathbb{Z}$, donc en \textit{définitive}
\[
  \boxed{\hat T_i \simeq \mathcal{O}(1)} \simeq \check F(\omega)
  \qquad \text{où } 0 \to F \to V(J) \to \mathcal{O}_\Sigma(1) \to 0 .
\]
On trouve donc
\[
  S\hat P_I \overset{\mathrm{déf}}{=} \prod_{/\hat M_{0,I}} \hat P_i
  \simeq \bigl(\check V^*(\mathcal{O}(1))^{I}\bigr)_{/\underset{\Sigma_I}{\hat M_{0,I}}}
  \simeq V^*(\mathcal{O}(1)) \wedge^{\mathbb{G}_m} \mathbb{G}_m^I
\]
donc
\[
  SP_I \simeq \bigl(V^*(\mathcal{O}(1)) \mid \underset{\Sigma_I^*}{M_{0,I}}\bigr) \wedge^{\mathbb{G}_m} \mathbb{G}_m^I
\]
d'où (sur $\mathbb{C}$ maintenant, et dans le contexte \uncertain{analytique})
\[
  S\mathcal{T}_{0,I} \simeq \Pi_1\bigl(V^*(\mathcal{O}(1)) \mid \Sigma_I^*\bigr) \wedge^{\mathbb{Z}} \mathbb{Z}^I
\]
\note{l'exposant $I$ de $\check V^*(\mathcal{O}(1))$ dans la première formule est lu avec doute ; le $\Pi$ de la dernière formule est récrit sur un premier tracé}

\section{Étude de $P\Sigma_J^*$ (et de $P\Sigma_n^*$ …) $S_0\mathcal{T}_{0,4}$, $S_0\mathcal{T}_{0,I}$ ($\mathrm{card}\, I = 4$)}
\note{titre de sa main, seul au haut du feuillet 18, qui sert de couverture à la suite}

\page{19}
\drawing{une page entière de dessin, au crayon : un cylindre vertical dont le cercle supérieur est marqué $V_{\mathbb{C}}^{\omega}$, le cercle inférieur $V_{\mathbb{C}}^{\omega'}$, et le cercle médian (l'équateur) $V^{\omega}$, chacun avec une flèche de sens de rotation. Sur le cercle supérieur, les points $\tilde P_{a_1}^{\omega}$, $\tilde P_s^{\omega}$, $\tilde P_a^{\omega}$, et l'arc $\tilde d_s^{\omega}$ (un premier libellé biffé) tracé en gras de $\tilde P_{a_1}^{\omega}$ à $\tilde P_a^{\omega}$ ; de $\tilde P_a^{\omega}$ descend en gras un segment vertical, libellé $\tilde c_r$ ou $\tilde c_a^{\omega}$, jusqu'au point $\tilde Q_a$ de l'équateur ; le long de la verticale en pointillé issue de $\tilde P_{a_1}^{\omega}$ on lit $\tilde c_{a_1}^{\omega} = \tilde c_{\omega^{-1}(r_1)}$, et plus bas $\tilde c_{r_1} = \tilde c_{a_1}^{\omega'}$. Sur l'équateur, quatre petits cercles fléchés autour de points marqués (des lacets autour des points enlevés), les points $\tilde Q_{a_1}$, $\tilde R_{r_1}$, $\tilde R_s$, $R_r$, $\tilde Q_a$, les arcs $\tilde b_{r_1}$, $\tilde b_r$ (en gras), et un demi-cercle gras $\tilde a_r$ autour de $\tilde R_s$ ; un triangle hachuré à l'intérieur porte les libellés $r$, $r'$, $r''$, $r_1$, $r_2$, plusieurs autres étant biffés. Sur le cercle inférieur, les points $\tilde P_{a_1}^{\omega'}$, $\tilde P_s^{\omega}$, $\tilde P_a^{\omega'}$, les libellés $(\tilde d_{r_1})$, $\tilde d_{r_1}$ ou $\tilde d_s^{\omega'}$, et l'arc fléché $\tilde b_r^{\omega'}$. Des lettres pâles $a_\omega$, $a$, $a_{\omega'}$ dans la marge gauche du cylindre.}

\note{légende écrite à gauche du dessin :}
\[
  \begin{aligned}
    & s \in S, \qquad a \in A, \qquad r = (s, a) \\
    & r_1 = (s, a_1) = \tau_1 r \\
    & r_2 = \omega^{-1}(r_1) = \tau_0 \tau_1 r = \tau_0 \tau_1 r \\
    & r' = \tau_0 r \\
    & r'' = \tau_1 r' = \tau_1 \tau_0 r = \omega(r)
  \end{aligned}
\]
\note{entre les lignes de $r_2$ et de $r'$, une ligne biffée : \struck{$\ldots = \omega(r) = \tau_0(\tau_1 r)$} ; dans la ligne de $r_2$, un symbole biffé et récrit après $\tau_0$. Le cercle entourant un chiffre en haut de la page n'est pas lu.}


\note{la suite annoncée par le titre du feuillet 18 (étude de $P\Sigma_J^*$ et de $S_0\mathcal{T}_{0,I}$ pour $\mathrm{card}\, I = 4$) se poursuit au-delà de ce lot ; le dessin du feuillet 19 n'en est que le début}

\end{document}
